% Generated draft: requires mathematical and source review.
\documentclass[11pt]{article}
\usepackage[margin=25mm]{geometry}
\usepackage{fontspec}
\setmainfont{Latin Modern Roman}
\usepackage{amsmath,amssymb,mathtools,unicode-math}
\setmathfont{Latin Modern Math}
\usepackage{enumitem,needspace,xurl,hyperref}
\setlength{\parindent}{0pt}
\setlength{\parskip}{5pt}
\setlength{\emergencystretch}{3em}
\newcommand{\N}{\mathbb{N}}
\newcommand{\Z}{\mathbb{Z}}
\newcommand{\Q}{\mathbb{Q}}
\newcommand{\R}{\mathbb{R}}
\newcommand{\C}{\mathbb{C}}
\newcommand{\F}{\mathbb{F}}
\newcommand{\A}{\mathbb{A}}
\newcommand{\PP}{\mathbb{P}}
\newcommand{\E}{\mathbb{E}}
\newcommand{\eps}{\varepsilon}
\newcommand{\dd}{\,\mathrm{d}}
\begin{document}
\section*{UnsolvedMath ID: 20000411}
\textbf{Source problem code:} catalogue entry
\subsection{Definitions and mathematical statement}
Let \(X/\Q\) be the Hilbert modular surface in the source discussion, and let \(\Q^{\mathrm{ab}}\) denote the maximal abelian extension of \(\Q\). For a smooth projective model and a prime \(\ell\), a divisor Tate class is a class in \(H^2_{\mathrm{\acute et}}(X_{\overline\Q},\Q_\ell(1))\) fixed by the Galois group of some finite extension of \(\Q\). A divisor represents such a class if its cycle class is that class, allowing rational or coefficient-field linear combinations as appropriate. Being defined over a field for a cohomology class means Galois invariance over that field; an actual cycle representative has a separate descent question.

Ramakrishnan's 2007 problem reports that the Tate classes arising from modular curves on \(X\) are defined over \(\Q^{\mathrm{ab}}\), while certain other classes are known to be algebraic and are defined over dihedral extensions, but not over \(\Q^{\mathrm{ab}}\). The task is to find specific algebraic representatives for those latter classes. This does not assert that every Hilbert modular surface has every type of class described. Non-congruence curves on \(X\) are suggested as possible representatives. Intersecting representatives with modular curves is proposed as a possible way to obtain interesting non-Heegner points, not as a proven conclusion. The surface's level, compactification/model, classes, and meaning of an explicit representative are not specified further.

The supplied later report offers a sufficient criterion rather than a representative. In the actual finite dihedral Galois quotient \(G\) governing a proposed divisor orbit, take an irreducible representation \(\rho\) of dimension greater than one. Over a characteristic-zero splitting coefficient field \(E\), the central projector on the corresponding isotypic component is
\[
 e_\rho=\frac{\dim\rho}{|G|}\sum_{\sigma\in G}\chi_\rho(\sigma^{-1})\sigma.
\]
In a divisor-class space carrying this Galois action and an invariant intersection form, put \(D=e_\rho[C]\) for a proposed curve \(C\) defined over the dihedral field. The reported character-weighted certificate is
\[
 \Theta_\rho(C)=\frac{\dim\rho}{|G|}
 \sum_{\sigma\in G}\chi_\rho(\sigma^{-1})(C\mathbin{\cdot}\sigma C)
 =[C]\mathbin{\cdot}D=D\mathbin{\cdot}D.
\]
The last equality uses the self-adjoint idempotent in this dihedral intersection setting. Nonzero \(\Theta_\rho(C)\) is sufficient to certify \(D\ne0\); the converse is not asserted. The report states that divisors defined over abelian extensions have zero \(\rho\)-projection, and that \(D\) has intersection degree zero with each abelian-defined modular-curve class. For a \(\Q\)-defined smooth projective modular curve \(M\), a well-defined restriction of this divisor combination therefore has degree zero. If its projected class is nonzero in \(\operatorname{Jac}(M)(\overline\Q)\otimes E\), the report concludes that the projected combination is non-torsion and cannot descend to \(\Q^{\mathrm{ab}}\). This conditional conclusion neither guarantees a nonzero restriction nor supplies an individual non-Heegner point. Compactification, proper intersection and restriction conventions, and the actual Galois module must be fixed before applying the criterion.
\subsection{Short English statement}
Find concrete curves representing the algebraic Tate classes on a Hilbert modular surface whose fields of definition are dihedral rather than abelian. The supplied report gives a character-projector and intersection certificate for proposed curves, but produces no explicit representative and no guaranteed non-Heegner point.
\subsection{Sources}
\begin{itemize}
\item Local UnsolvedMath registry, permanent ID \(20000411\), source code \texttt{catalogue entry}; AIM algebraic number theory corpus, \texttt{aim-algebraic-number-theory-notes.json}, notes[62]. Attribution: Ulam AI / UnsolvedMath Contributors, CC BY 4.0.
\item \emph{Problems from the AIM Tate Conjecture Workshop}, July 23--27, 2007, transcribed by Christopher Lyons, Problem 2 and both comments attributed to D. Ramakrishnan: \url{https://aimath.org/WWN/tateconjecture/tateconjecture.pdf}. These identification details are supplied by the verification field; no new source inspection is claimed. Original literature and remarks fields are empty.
\item Registry dated August 14, 2026: statement exact, status \texttt{partially\_solved}, classification \texttt{PARTIAL-PROGRESS}, suggested difficulty L4; attempt 1 records problem status open, result reduction, rigor proved, novelty confidence low. The projector/intersection/Jacobian procedure is a machine-generated research aid, not peer reviewed, requiring independent verification before citation; no primary bibliographic reference for it is supplied.
\end{itemize}
\end{document}
